\documentclass[10pt,a4paper]{amsart}
\usepackage[margin=19mm]{geometry}
\usepackage{amsmath,amssymb,mathtools}
\usepackage[scr=rsfs,cal=euler]{mathalfa}
\usepackage{xfrac}
\usepackage[unicode,hidelinks,bookmarks=false]{hyperref}
\newcommand{\ie}{i.e.\ }
\newcommand{\cf}{cf.\ }
\newcommand{\resp}{respectively\ }
\newcommand{\Subsection}{\subsection}
\providecommand{\nobreakdash}{\nobreak\hbox{-}}
\setlength{\emergencystretch}{2em}
\multlinegap0pt
\usepackage{bm}
\usepackage{bbm}
\usepackage{dsfont}
\usepackage{xspace}
\usepackage{cancel}
\usepackage{mathtools}
\usepackage{amsthm}
\makeatletter
\@ifundefined{thm}{\newtheorem{thm}{Thm}}{}
\@ifundefined{lem}{\newtheorem{lem}[thm]{Lemma}}{}
\makeatother
\title[Minimal dimension equivariant embeddings of real and complex]{Minimal dimension equivariant embeddings of real and complex flag manifolds into Euclidean spaces}
\author{Zhongzi Wang, Hang Yin}
\date{Source: arXiv:2605.17337v1 (CC BY 4.0)}
\begin{document}
\maketitle

\noindent\emph{Source and licence.} Minimal dimension equivariant embeddings of real and complex flag manifolds into Euclidean spaces. Version arXiv:2605.17337v1 (CC BY 4.0), distributed under the Creative Commons Attribution 4.0 International licence, as recorded in the arXiv licence metadata for this exact version and shown on the abstract page https://arxiv.org/abs/2605.17337v1. Full rights notice: licenses/arxiv-2605.17337-CC-BY-4.0.txt.\par\smallskip

\noindent\textit{Editorial scope.} Counted unit: the Lem whose source label is \texttt{first type1} in the article (source0.tex lines 834--840, source label \texttt{first type1}), delivered together with the article's own complete original proof of it (source0.tex lines 841--876, one \texttt{proof} environment of 36 lines). No mathematics is written, repaired or re-derived: statement and proof are the article's own text line for line. Required context retained uncounted: case 1 (lines 804--806); case 2 (lines 808--813); case 2.3 (lines 824--826); case 3 (lines 828--830). Every definition, display and label of those lines is retained unchanged. Omitted from this package and not redistributed: 1--803, 807--807, 814--823, 827--827, 831--833, 877--1356. Nothing inside a retained excerpt is omitted or altered: every retained body is delivered exactly as the article's lines are written, so no omission receipt of any kind accompanies this package. Citations: the retained excerpts carry no citation command at all, so no bibliography entry of the article is transcribed and no reference is redistributed. Reference adaptation: none was needed and none was made; every reference inside the delivered unit resolves inside it, so no unresolved reference or citation is produced. Printed slips kept literal: no printed word, symbol, hypothesis or number of the counted statement or of its proof was altered. Layout and class: the article's own class is \texttt{amsart} with packages including amscd, amsthm, amsfonts, amssymb, amsmath, esint, tikz-cd, hyperref, fullpage, xy, mathtools, color; the wrapper is the lane's pinned \texttt{amsart} editorial wrapper at 10pt on A4, which restates the theorem-like environments the retained text needs and adds the article's own macro definitions that the retained text needs (none), with extra wrapper packages (bm, bbm, dsfont, xspace, cancel, mathtools). The delivered numbering is the unit's own. Assets: no retained span contains a figure environment, a graphics-inclusion command, a PDF-page inclusion, a table or a TikZ picture; no image file is copied into this package, the package declares no asset dependency and no image file is redistributed with it. Licence: version arXiv:2605.17337v1 of this article is distributed under the Creative Commons Attribution 4.0 International licence, as recorded in the arXiv licence metadata for this exact version and shown on the abstract page https://arxiv.org/abs/2605.17337v1. A full rights notice is delivered at licenses/arxiv-2605.17337-CC-BY-4.0.txt. Characters: every retained excerpt is pure ASCII, so the delivered engine, XeLaTeX with its default Latin Modern text font, reports no missing character.
		\begin{equation}\label{case 1}
		    \dim V(\mu)=n(2n+3).
		\end{equation}

		\begin{equation}\label{case 2}
		    \begin{aligned}
		    \dim V(\mu)&=\prod_{1\leqslant i\leqslant k}\dfrac{n+1-i}{k+1-i}\prod_{1\leqslant i\leqslant k}\dfrac{2n+1-i-k}{n+1-i}\prod_{1\leqslant i\leqslant j\leqslant k}\dfrac{2n+3-i-j}{2n+1-i-j}\\
            &=\prod_{1\leqslant i\leqslant k}\dfrac{2n+1-i-k}{k+1-i}\prod_{1\leqslant i\leqslant j\leqslant k}\dfrac{2n+3-i-j}{2n+1-i-j}.
		\end{aligned}
		\end{equation}

        \begin{equation}\label{case 2.3}
            \dim(V(1,1,1,,0,\dots,0))=\dfrac{n(2n+1)(2n-1)}{3}.
        \end{equation}

    \begin{equation}\label{case 3}
         \dim V(1,1,\dots,1)=\prod_{1\leqslant i\leqslant j\leqslant n}\dfrac{2n+3-i-j}{2n+1-i-j}.
    \end{equation}

\begin{lem}\label{first type1}
	%When $k\geqslant 3$ and $n\geqslant 4$, we have
	%\[\dim V(\underbrace{1,\ldots,1}_k,0,\ldots,0)\geqslant\dim V(2,0,\ldots,0).\]
    When $k\geqslant 3$, we have
	\[\dim V(\underbrace{1,\ldots,1}_k,0,\ldots,0)>\dim V(2,0,\ldots,0).\]
    
\end{lem}

\begin{proof}
	Denote $\mu=(2,0,\ldots,0)$ and $\lambda_k=(\underbrace{1,\ldots,1}_k,0,\ldots,0)$. First we prove that when $k\leqslant n-1$, we have $\dim V(\lambda_{k+1})>\dim V(\lambda_k)$. 

Note that by equation ~\eqref{case 1}, when $k<n-1$, 
    \begin{equation}\label{quotient of V(mu)}
    \begin{aligned}
        \dfrac{\dim(V(\lambda_{k+1}))}{\dim(V(\lambda_k))}&=\dfrac{2(n-k)(2n-1-2k)}{(2n-k)(k+1)}\cdot \dfrac{(2n+1-k)(2n-k)}{2(n-k)(2n-2k-1)}\\
        &=\dfrac{2n+1-k}{k+1}\geqslant \dfrac{k+2}{k+1}>1
    \end{aligned}
    \end{equation}
    And by equations ~\eqref{case 2} and ~\eqref{case 3},
    \[\begin{aligned}
    \dfrac{\dim V(\lambda_n)}{\dim V(\lambda_{n-1})}&=\dfrac{n+1}{n}>1
    \end{aligned}
    \]
%\blue{
%We have
%\[\frac{\dim V(\lambda_{k+1})}{\dim V(\lambda_k)}=\frac{2n-k+1}{k+1}\geqslant\frac{k+3}{k+1}>1.\]

%}
  %  Obviously when $k$ increases, the last quantity in ~\eqref{quotient of V(mu)} decreases, and when $k=n-1$, it equals to $\dfrac{n+2}{3}>1$ for every $n>1$, so $\dim(V(\lambda_{k+1}))>\dim(V(\lambda_{k}))$ for each $1\leqslant k\leqslant n-1$. 

Thus to finish the proof of the lemma, we only need to show that for $n\geqslant 3$,
\begin{equation}
    \dim(V(\lambda_3))> \dim(V(\mu))
\end{equation}
This is clear from equations ~\eqref{case 1} and ~\eqref{case 2.3}, since when $n\geq 3$,
\[
\dfrac{n(2n+1)(2n-1)}{3}-n(2n+3)=n\cdot \dfrac{4n^2-6n-10}{3}>0
\]

%\blue{
%It is clear for all $n\geqslant 3$,
%\[\dim V(\lambda_3)>\dim V(\mu).\]
%}
\end{proof}

\end{document}
